\scanpage{21}

\unsure{pondent} respectivement aux $\tau$ qui \unsure{avec} $1$ \unsure{engendrent} \unsure{l'un} réseau des entiers alg. de $\mathbb{Q}(i)$ (soit $\mathbb{Z}+i\mathbb{Z}$), et le réseau des entiers de $\mathbb{Q}(j)$, formé des nbs de la forme $aj+b\bar{j}$ ($j=\dfrac{-1+i\sqrt{3}}{2}=\exp\dfrac{2i\pi}{3}$); la ramification \unsure{correspondante} est $2$ et $3$.

On a un isom. canonique (analytique) \inserted{\unsure{entre}}
\[
\tag{29}
\unsure{\underline{D}|D}/\mathrm{GP}(1,\mathbb{Z})^{+}\simeq\mathbb{P}^{1\,\mathrm{an}}\setminus\{0\}=Y,
\]
\unsure{caractérisé} par la condition que l'orbite critique de ramification $3$ \unsure{corresponde} au point $\infty$, et celle de ramification $2$ au point $1$. On trouve ainsi un isom
(posant $\unsure{\mathcal{G}|\mathfrak{G}}=\mathrm{GP}(1,\mathbb{Z})$, $G=\mathrm{GP}(1,\mathbb{Z})^{+}=\unsure{\mathrm{SL}}(2,\mathbb{Z})'$)

(30)
\diagram{
\begin{tikzcd}[ampersand replacement=\&,column sep=large]
\pi_1(\unsure{\underline{D}|D},\unsure{\mathcal{G}|\mathfrak{G}},x)
  \arrow[r,"\sim"] \arrow[d,equal]
\& \pi_1^{\unsure{*|\pm}}(Y,\underline{d};y) \\
\unsure{\mathcal{G}|\mathfrak{G}} \&
\end{tikzcd}
}
\margin{$(x\in\unsure{\underline{D}|D}\setminus\unsure{D^{!}|D'})$}

que je voudrais expliciter, en identifiant \deleted{les} $\tau'_0,\tau'_1,\tau'_\infty$ comme \unsure{éléments} de $\unsure{\mathrm{GL}}^{\unsure{\pm'|\pm}}(2,\mathbb{Z})'$.

\margin{\unsure{0}}
\margin{20}

\scanpage{22}
\mspage{299}

\uline{39} \deleted{\unsure{Quelques}}\inserted{Les} conditions de congruence
pour l'opération extérieure de $\Gamma_{\mathbb Q}$ sur $\widehat{\Pi}_{0,3}$.

Considérons le diagramme

\[
\tag{1}
\diagram{\begin{tikzcd}[ampersand replacement=\&,column sep=large,row sep=small]
1 \arrow[r] \& \unsure{\widehat{\mathcal T}_{1,1}^{+}} \arrow[r] \arrow[d,phantom,"\simeq" description,sloped] \& N_{1,1} \arrow[r] \arrow[dd] \& \Gamma_{\mathbb Q} \arrow[r] \arrow[dd,"\chi"] \& 1 \\
\& \unsure{\mathrm{SL}}(2,\mathbb Z)^{\wedge} \arrow[d] \& \& \& \\
1 \arrow[r] \& \unsure{\mathrm{SL}}(2,\widehat{\mathbb Z}) \arrow[r] \& \unsure{\mathrm{GL}}(2,\widehat{\mathbb Z}) \arrow[r,"\det"] \& \widehat{\mathbb Z}^{*} \arrow[r] \& 1
\end{tikzcd}}
\]

Il montre que l'opération extérieure de
$\Gamma_{\mathbb Q}$ sur $\unsure{\widehat{\mathcal T}_{1,1}^{+}}$
\inserted{$\simeq \unsure{\mathrm{SL}}(2,\mathbb Z)^{\wedge}$}
passe \unsure{en|au} quotient en une
action extérieure sur $\unsure{\mathrm{SL}}(2,\widehat{\mathbb Z})$, et cette
action extérieure, de plus, peut s'expliciter
entièrement à l'aide du caractère
cyclotomique
\[
\tag{2}
\chi : \Gamma_{\mathbb Q}\longrightarrow \widehat{\mathbb Z}^{*}.
\]

Pour préciser ce dernier point, notons
que la deuxième ligne de (1) donne
un homomorphisme
\[
\tag{3}
\widehat{\mathbb Z}^{*}\longrightarrow
\mathrm{Autext}(\unsure{\mathrm{SL}}(2,\widehat{\mathbb Z})),
\]
et l'hom canonique
\[
\tag{4}
\Gamma_{\mathbb Q}\longrightarrow
\deleted{\unsure{\unsure{\mathrm{SL}}(2,}}
\mathrm{Autext}(\unsure{\mathrm{SL}}(2,\widehat{\mathbb Z}))
\]
\margin{21}

\scanpage{23}

n'est autre que le composé de (2) et (3).
Par ailleurs, l'hom. (4) se factorise
en
\[
\tag{5}
\widehat{\mathbb Z}^{*}/\widehat{\mathbb Z}^{*2}
\longrightarrow \mathrm{Autext}(\unsure{\mathrm{SL}}(2,\widehat{\mathbb Z}))
\]

(\unsure{provient} du fait que les \unsure{éléments|élts} de
$\unsure{\mathrm{GL}}(2,\widehat{\mathbb Z})$ qui sont dans le centre
$\widehat{\mathbb Z}^{*}$ opèrent trivialement
\deleted{\unsure{du point}}\inserted{sur} $\unsure{\mathrm{SL}}(2,\widehat{\mathbb Z})$,
et que le composé
\[
\widehat{\mathbb Z}^{*}\xrightarrow{\text{(centre)}}
\unsure{\mathrm{GL}}(2,\widehat{\mathbb Z})\xrightarrow{\det}
\unsure{\widehat{\mathbb Z}|\mathbb Z}^{\unsure{*|2}}
\]
n'est autre que $\lambda\longmapsto\lambda^{\unsure{2}}$). Plus
précisément, on voit que pour que l'
\unsure{autom} de $\unsure{\mathrm{SL}}(2,\widehat{\mathbb Z})$ induit par un
$u\in\unsure{\mathrm{GL}}(\unsure{1|2},\widehat{\mathbb Z})$ soit intérieur,
\unsure{i.d.|i.e.} de la forme
$\mathrm{int}\,\unsure{(u_0)|u_0}$ avec $u_0\in\unsure{\mathrm{SL}}(2,\widehat{\mathbb Z})$, il f. et s.
qu'il existe $u_0\in\unsure{\mathrm{SL}}(\unsure{2|1},\widehat{\mathbb Z})$ tel que
$u^{\unsure{\deleted{0}|\circ}}u_0^{-1}$ commute à $\unsure{\mathrm{SL}}(2,\widehat{\mathbb Z})$, et cela signifie
que c'est un scalaire -- i.e.
$u\in\unsure{\mathrm{SL}}(2,\widehat{\mathbb Z})\cdot\widehat{\mathbb Z}^{*}$\unsure{,|.}
En d'autres termes
l'hom (5) est \uline{injectif}, et en conséquence,
\margin{22}

\scanpage{24}
\mspage{300}

le noyau de (4) n'est autre que le
noyau du composé
\[
\tag{6}
\Gamma_{\mathbb Q}\xrightarrow{\chi}\widehat{\mathbb Z}^{*}
\longrightarrow\widehat{\mathbb Z}^{*}/\widehat{\mathbb Z}^{*2}
\]
i.e. (par la théorie du corps de classes pour
$\mathbb Q$) l'opération \inserted{ext} de $\Gamma_{\mathbb Q}$ sur
$\unsure{\mathrm{SL}}(\unsure{1|2},\widehat{\mathbb Z})$ se
factorise par une opération \inserted{ext} de
$\unsure{(\Gamma_{\mathbb Q\,\mathrm{ab}})_2|(\Gamma_{\mathbb Q}^{\mathrm{ab}})_2}$ -- le
\unsure{quotient} ``multiquadratique'' de $\Gamma_{\mathbb Q}$ -- et cette
dernière est fidèle.

\margin{\uline{NB} L'op. ext. de $\Gamma_{\mathbb Q}$ sur
$\unsure{\mathrm{SL}}(2,\mathbb Z/\unsure{n|m}\mathbb Z)$ se fait par
l'\unsure{intermédiaire} de
$\unsure{((\mathbb Z/\unsure{n|m}\mathbb Z)^{*})_2|(\mathbb Z/\unsure{n|m}\mathbb Z)_2^{*}}$}

Ainsi, si on identifie $\Gamma_{\mathbb Q}$ \deleted{\unsure{par (4)}}
à un sous-groupe de $\mathrm{Autext}\,\unsure{\mathrm{SL}}(2,\widehat{\mathbb Z})^{\wedge}$,
ce sous-groupe \inserted{$\unsure{\Gamma}$} est astreint à des
conditions assez fortes : \deleted{\illegible{1}} \unsure{invarier}
le sous-groupe distingué noyau de
\[
\tag{7}
\unsure{\mathrm{SL}}(2,\mathbb Z)^{\wedge}\longrightarrow
\unsure{\mathrm{SL}}(2,\widehat{\mathbb Z})
\]
et induire sur $\unsure{\mathrm{SL}}(2,\widehat{\mathbb Z})$
\deleted{\unsure{l'}}\inserted{une} opération extérieure
provenant d'un hom $\Gamma\longrightarrow
\widehat{\mathbb Z}^{*}/\widehat{\mathbb Z}^{*2}$.
On peut d'ailleurs préciser cet hom.
ainsi. Considérons dans $\unsure{\mathrm{SL}}(2,\widehat{\mathbb Z})$ le
``\unsure{sous-groupe-lacets|sous-groupe-lacet}'' provenant du point \unsure{à|:}
\margin{23}

\scanpage{25}

l'infini de $M_{1,1\,\mathbb Q}$, il correspond
à une classe d'hom. ext.
\[
(8)\qquad
\overline T=\overline T_{\infty}(\overline{\mathbb Q}_{0})
\longrightarrow
\unsure{\mathrm{SL}}(2,\overset{\unsure{\deleted{\wedge}|\wedge}}{\mathbb Z})^{\wedge}
\]
dont l'image \inserted{dans $\unsure{\mathrm{SL}}(2,\widehat{\mathbb Z})$}
sauf erreur, est (à conj. près) un sous-groupe ``unipotent''
\[
(9)\qquad
\left(\begin{array}{cc}1&0\\ n&1\end{array}\right),
\qquad n\in\widehat{\mathbb Z}.
\]

Cet hom extérieur est respecté par l'opération de $\Gamma$, d'où une opération
bien déterminée de $\Gamma$ sur \deleted{\illegible{1}} $\overline T$
i.e. un hom
\deleted{$\Gamma\longrightarrow\operatorname{Aut}(\overline T)\simeq\widehat{\mathbb Z}^{*}$}
\deleted{(qui \unsure{sera} \illegible{1}, justement $\unsure{\chi}$, le}
\[
(10)\qquad
\Gamma\xrightarrow{\chi}\operatorname{Aut}(\overline T)
\simeq\widehat{\mathbb Z}^{*},
\]
qui \unsure{sera} justement le \unsure{caractère|caractre} cyclotomique.
De façon plus précise, pour \unsure{un|Un} automorphisme extérieur \inserted{$u$} de
$\unsure{\mathrm{SL}}(2,\mathbb Z)^{\wedge}$ (ou seulement de $\unsure{\mathrm{SL}}(2,\widehat{\mathbb Z})$)
qui invarie \unsure{la|l'} \deleted{\unsure{im.}} \unsure{classe} de \unsure{conjugais.|conjugais}
de l'image de (8) [ou \unsure{seulement|seulmt} du composé
$\overline T\unsure{\longrightarrow}\unsure{\mathrm{SL}}(2,\widehat{\mathbb Z})$],
\inserted{on vérifie qu'}il existe un
$\chi(u)\in\operatorname{Aut}(\overline T)\simeq\widehat{\mathbb Z}^{*}$ unique,
tel que l'on ait commutativité dans
\margin{faire la vérification}
\margin{24}

\scanpage{26}
\mspage{301}

\[
(11)\qquad
\diagram{\begin{tikzcd}[ampersand replacement=\&]
\overline T \arrow[r, "\unsure{\kappa|k}"] \arrow[d, "\chi(u)"]
 \& \unsure{\mathrm{SL}}(2,\mathbb Z)^{\wedge} \arrow[d, "u"] \\
\overline T \arrow[r]
 \& \unsure{\mathrm{SL}}(2,\mathbb Z)^{\wedge}
\end{tikzcd}}
\]
--- ou \unsure{même}, seulement dans
\[
(\unsure{12|11})\qquad
\diagram{\begin{tikzcd}[ampersand replacement=\&]
\overline T \arrow[r, "\unsure{\kappa|k}"] \arrow[d, "\chi(u)"]
 \& \unsure{\mathrm{SL}}(2,\widehat{\mathbb Z}) \arrow[d, "\dot u"] \\
\overline T \arrow[r]
 \& \unsure{\mathrm{SL}}(2,\widehat{\mathbb Z})
\end{tikzcd}}
\quad ---
\]
d'où un \unsure{caractère|caractre} \inserted{$\chi$} (10) sur le
groupe de tous les autom. \inserted{ext.} $u$ de
$\unsure{\mathrm{SL}}(2,\widehat{\mathbb Z})$ qui ont la propriété \unsure{envisagée}.

Ceci dit, $\Gamma$ est contenu dans ce sous-groupe \inserted{$\unsure{\Gamma}$},
et de plus dans le sous-groupe des autom. ext. qui \uline{de plus}
préservent le noyau de (7), et qui opèrent \inserted{(ext.!)} dans le quotient
$\unsure{\mathrm{SL}}(2,\widehat{\mathbb Z})$ suivant le composé
\[
\Gamma\unsure{:}\xrightarrow{\chi}\widehat{\mathbb Z}^{*}
\longrightarrow\widehat{\mathbb Z}^{*}/\widehat{\mathbb Z}^{*2}
\hookrightarrow\operatorname{Autext}_{\unsure{\mathrm{lac}!|\mathrm{loc}!}}
(\unsure{\mathrm{SL}}(2,\widehat{\mathbb Z})).
\]

\deleted{L'on} D'ailleurs notons que le centre $\{\pm1\}$
de $\unsure{\mathrm{SL}}(2,\mathbb Z)^{\wedge}$ n'étant pas trivial, l'op. ext.
\margin{25}

\scanpage{27}

de $\Gamma_{\mathbb Q}$ sur
$\unsure{\widehat{\mathcal T}_{1,1}^{+}|\widehat{\tau}_{1,1}^{+}}$
ne permet pas à priori de reconstituer (à iso unique près)
l'ext. $N_{1,1}$ de $\Gamma_{\mathbb Q}$ par
$\unsure{\widehat{\mathcal T}_{1,1}^{+}|\widehat{\tau}_{1,1}^{+}}$
de (1). Elle donne seulement l'extension \deleted{\illegible{1}}
$N'_{1,1}$ de $\Gamma_{\mathbb Q}$ par
\[
\unsure{\widehat{\mathcal T}_{1,1}^{+\prime}|\widehat{\tau}_{1,1}^{+\prime}}
\simeq\unsure{\mathrm{SL}}(2,\mathbb Z)^{\wedge}/\{\pm1\}
\inserted{\overset{\unsure{\mathrm{dfn}|\mathrm{déf}}}{=}
\unsure{\mathrm{SL}}(2,\widehat{\mathbb Z})'}
\]
(par image inverse de l'extension

\diagram{[(13) : suite d'objets et de flèches vers la droite,
$1\longrightarrow\unsure{\unsure{\mathrm{SL}}(2,\widehat{\mathbb Z})'}
\longrightarrow
\operatorname{Aut}\deleted{\unsure{\mathrm{ext}}}_{\unsure{\mathrm{lac}|\mathrm{loc}}}
\bigl(\unsure{\mathrm{SL}}(2,\unsure{\mathbb Z|\dot{\mathbb Z}})^{\wedge}\bigr)
\longrightarrow
\unsure{\operatorname{Autext}_{\mathrm{lac}}|\operatorname{Autext}_{\mathrm{loc}}}
\unsure{\unsure{\mathrm{SL}}(2,\mathbb Z)^{\wedge}}$.
Le dernier terme est serré contre le bord droit; sa parenthèse et ses éventuelles
primes ne sont pas lisibles. Une petite flèche descendante à droite mène à $1$.]}

Cette extension \inserted{$N'_{1,1}$} a été aussi interprétée
comme l'extension $M_{0,3}\simeq N'_{0,4}$ de
$\Gamma_{\mathbb Q}$ par
$\unsure{\widehat{\mathfrak S}_{0,3}|\widehat{\mathfrak G}_{0,3}}
\simeq\unsure{\widehat{\mathcal T}_{0,4}'|\widehat{\tau}_{0,4}'}$
(groupe qui est lui-$\overline{\mathrm m}$ ext de
$\unsure{\mathfrak S_3|\mathfrak G_3}\simeq\unsure{\mathrm{SL}}(2,\mathbb Z/2\mathbb Z)$
par $\widehat{\Pi}_{0,3}$).
Nous savons déjà que l'opération \inserted{ext} de $\Gamma_{\mathbb Q}$
sur le noyau $\widehat{\Pi}_{0,3}$ de
$\unsure{\mathrm{SL}}(2,\overset{\unsure{\deleted{\wedge}|\wedge}}{\mathbb Z})^{\wedge\prime}
\longrightarrow\unsure{\mathrm{SL}}(2,\unsure{\mathbb Z/2|\mathbb Z_{2}})$
est fidèle, et commute à l'op. ext. de
$\unsure{\mathfrak S_3|\mathfrak G_3}\simeq\unsure{\mathrm{SL}}(2,\mathbb Z/2\mathbb Z)
\simeq\unsure{\mathrm{SL}}(2,\mathbb Z)^{\wedge\prime}/\widehat{\Pi}_{0,3}$.
Mais cette commutation peut mainte-
\margin{26}

\scanpage{28}
\mspage{302}
nant s'interpréter comme conséquence du fait que l'opération ext. de $\Gamma_{\mathbb Q}$ sur $\unsure{\mathfrak S_3|\mathfrak G_3}\simeq \unsure{\mathrm{SL}}(2,\mathbb Z/2\mathbb Z)$ est triviale, elle-$\overline{\mathrm m}$ \inserted{(dans l'optique présente)} conséquence du fait qu'elle \unsure{se fait} par l'intermédiaire d'une op. (fidèle) de $(\mathbb Z/n\mathbb Z)^*_{\unsure{/2|2}}$, or pour $n=2$ $(\mathbb Z/n\mathbb Z)^*_2=1$ (relation qui est par contre fausse pour \deleted{\unsure{les}} \inserted{\unsure{autres|d'autres}} \unsure{valeurs}, $n\neq\deleted{\unsure{2,}}1,2$ i.e. $n\geq3$). Cette relation de commutation, de nature assez triviale, apparaît donc à présent comme un cas très particulier des relations congruentielles nettement plus subtiles. On peut \inserted{essayer de} les exprimer \deleted{\unsure{de}} \deleted{façon} directement en termes de l'opération ext. de $\Gamma_{\mathbb Q}$ sur $\widehat\Pi_{0,3}$, en introduisant les noyaux des \unsure{hom.} $\unsure{\widehat{\mathcal T}_{1,1}^{+\prime}|\widehat\tau_{1,1}^{+\prime}}\longrightarrow \unsure{\mathrm{SL}}(2,2m)$, qui sont des sous-groupes \deleted{$\Pi(2m)$} ouverts :
\[
(14)\qquad \Pi(2m)\subset\widehat\Pi_{0,3}
\]
distingués et invariants par $\unsure{\mathfrak S_3|\mathfrak G_3}$, noyaux
\margin{27}

\scanpage{29}
d'homomorphismes
\[
(15)\qquad \varphi_{2m}:\widehat\Pi_{0,3}\longrightarrow\unsure{\mathrm{SL}}(2,\mathbb Z/2m)'
\]
qu'il \unsure{ne|en} doit pas être bien difficile d'expliciter, \deleted{\unsure{en termes}} \inserted{par ``réduction mod $2m$''} d'un \unsure{hom.}
\[
(16)\qquad \varphi_{\infty}:\widehat\Pi_{0,3}\longrightarrow\unsure{\mathrm{SL}}(2,\widehat{\mathbb Z})'
\]
dont il a été question déjà précédemment -- \unsure{hom.} dont l'image est le sous-groupe noyau de $\unsure{\mathrm{SL}}(2,\widehat{\mathbb Z})'\longrightarrow\unsure{\mathrm{SL}}(2,\mathbb Z/2\mathbb Z)$, i.e. on a une suite exacte
\[
(17)\qquad \deleted{\illegible{3}}\inserted{1\longrightarrow\Pi(\infty)\longrightarrow}\widehat\Pi_{0,3}\longrightarrow\unsure{\mathrm{SL}}(2,\widehat{\mathbb Z})'\longrightarrow\unsure{\mathrm{SL}}(2,\mathbb Z/2\mathbb Z)\longrightarrow1
\]
\[
(18)\qquad 1\longrightarrow\Pi(2m)\longrightarrow\widehat\Pi_{0,3}\longrightarrow\unsure{\mathrm{SL}}(2,\mathbb Z/2m\mathbb Z)'\longrightarrow\unsure{\mathrm{SL}}(2,\mathbb Z/2\mathbb Z)\longrightarrow1.
\]
Ceci posé, \unsure{ces|les} suites exactes permettent de façon naturelle de définir une opération ext. de $\widehat{\mathbb Z}^{*}/\widehat{\mathbb Z}^{*2}$ sur $\widehat\Pi_{0,3}/\Pi(\infty)$, ou encore des $(\mathbb Z/2m\mathbb Z)^*_2$ sur les $\widehat\Pi_{0,3}/\Pi(2m)$, et la ``condition de \unsure{congruence|congruences}'' s'exprime en disant que le caractère cyclotomique $\chi$ de $\Gamma_{\mathbb Q}$ (\inserted{\unsure{qu'on peut considérer comme}} déduit de l'opération ext. de
\margin{28}

\scanpage{30}
\mspage{303}
$\Gamma_{\mathbb Q}$ sur \unsure{les} groupes \unsure{à lacets} $\widehat\Pi_{0,3}$), qui donne \unsure{un|en} composé
\[
(19)\qquad \Gamma_{\mathbb Q}\xrightarrow{\chi}\widehat{\mathbb Z}^{*}\xrightarrow{\mathrm{can.}}\widehat{\mathbb Z}^{*}/\widehat{\mathbb Z}^{*2},
\]
\unsure{de|donc} une op. ext. de $\Gamma_{\mathbb Q}$ sur $\widehat\Pi_{0,3}/\Pi(\infty)$, est \uline{compatible} avec l'op. ext. de $\Gamma_{\mathbb Q}$ sur $\widehat\Pi_{0,3}$ lui-$\overline{\mathrm m}$.

Pour chaque entier $m$, cette condition s'explicite par le fait que $\Pi(2m)$ est invariant par l'op. ext. de $\Gamma_{\mathbb Q}$, et que l'opération ext. induite \unsure{de|sur} $\widehat\Pi_{0,3}/\Pi(2m)$ est celle déduite du caractère composé de (19) avec
\[
\widehat{\mathbb Z}^{*}/\widehat{\mathbb Z}^{*2}\longrightarrow(\mathbb Z/\unsure{2m|m}\mathbb Z)^*_2:
\]
\[
(20)\qquad \Gamma_{\mathbb Q}\longrightarrow(\mathbb Z/m\mathbb Z)^*_2.
\]
On devrait regarder si cette formulation est vraiment \uline{équivalente} à celle des conditions de \unsure{congruences|congruence} dans $\unsure{\widehat{\mathcal T}_{1,1}^{+\prime}|\widehat\tau_{1,1}^{+\prime}}$ \unsure{tout} entier, \deleted{(\unsure{quitte à y} \illegible{1}} \inserted{et non seulement une} conséquence; \unsure{exercice|exercices} ! Par contre, \unsure{le|la}
\margin{29}
